\section{Positions of centres at the current stage}

The literal obligation in Bierstone--Milman \cite{bierstone-milman97-invent-math},
Theorem~1.6(1), is a statement about each current strict transform: a centre is
contained in its singular locus, or the current transform is smooth and the
centre is contained in its intersection with the accumulated divisor. The
results in this section distinguish the part forced by a marked-ideal tower
from the additional producer needed for that literal obligation.

\begin{lemma}[Transport away from the accumulated divisor]
\label{hir:stage-cosupport-off-boundary}
\uses{hir:geo-transform,hir:admissible-calculus}
\lean{Hironaka.HironakaStageCenters.preimage_boundary_subset_boundary_geoTransform,Hironaka.HironakaStageCenters.notMem_boundary_geoTransform,Hironaka.HironakaStageCenters.notMem_initial_boundary_of_notMem_boundary,Hironaka.HironakaStageCenters.exceptionalLocus_subset_boundary,Hironaka.HironakaStageCenters.mem_cosupp_geoTransformAlong_iff_of_notMem_boundary,Hironaka.HironakaStageCenters.mem_cosupp_iff_mem_strictTransform_of_notMem_boundary,Hironaka.HironakaStageCenters.cosupp_subset_strictTransform_union_boundary}
\leanok
Let $\mathfrak m=(X,N,E_0,\mathfrak a,d)$ have closed smooth $N$ and coherent
$\mathfrak a$, and let $T$ be a multiple test blowup. At a point $y$ outside
the accumulated divisor $|E_T|$, membership in the current cosupport is
equivalent to membership of $T(y)$ in the original cosupport. In particular,
for $\mathfrak m=(X,X,E_0,\mathcal I_W,1)$ with $W$ closed analytic, this is
equivalent to $y\in W_T$, where $W_T$ is the actual strict transform. Thus
$\operatorname{cosupp}(\mathfrak m_T)\subset W_T\cup|E_T|$.
\end{lemma}
\begin{proof}
Induct on the recorded tower. There is nothing to transform at the initial
stage. Write a final step as $\pi:Z\to Y$ with centre $C$. The new divisor
contains $\pi^{-1}C$, so $z\notin|E_{T\pi}|$ implies $\pi(z)\notin C$.
It also implies $\pi(z)\notin|E_T|$: if $\pi(z)\in H\in E_T$ and
$\pi(z)\notin C$, then $z\in\pi^{-1}(H\setminus C)$ and hence belongs
to the recorded strict transform of $H$. Off $\pi^{-1}C$ the controlled
transform agrees, locally up to a unit, with the pullback of the old ideal;
the blowup is a local biholomorphism. Consequently its cosupport is the
preimage of the old cosupport there. Apply the induction hypothesis.

For the final specialization, the original cosupport is exactly $W$. The
exceptional locus of the tower projection is contained in $|E_T|$, by
induction using the exceptional-locus composition formula. Away from that
locus the strict transform of a closed set is its inverse image. This proves
the equivalence. For any point of the cosupport, either it belongs to $|E_T|$
or the equivalence puts it in $W_T$.
\end{proof}

\begin{lemma}[Regularity and singularity at an unramified point]
\label{hir:stage-regularity-unramified}
\uses{ana:strict-transform-composition}
\lean{Hironaka.HironakaStageCenters.regularAtDim_strictTransform_iff_of_notMem_exceptionalLocus}
\leanok
Let $f:Y\to X$ be holomorphic between complex manifolds of the same dimension,
$W\subset X$ closed, and $y\notin\operatorname{Exc}f$. For every $d$,
\[
 \operatorname{Reg}_d(W_f,y)\quad\Longleftrightarrow\quad
 \operatorname{Reg}_d(W,f(y)).
\]
\end{lemma}
\begin{proof}
The holomorphic inverse function theorem gives open neighbourhoods of $y$
and $f(y)$ and a holomorphic inverse section $s$ through $y$. Restricting $s$
to its open domain gives an actual holomorphic map of complex manifolds.
The chain rule applied to $f\circ s=\mathrm{id}$ shows that both
differentials are bijective. The inverse section therefore transports
flattening coordinates in both directions; this identifies regularity of
$f^{-1}W$ at $y$ with regularity of $W$ at $f(y)$.
The exceptional locus is closed, and on its open complement the strict
transform of closed $W$ has the same germ as $f^{-1}W$. Replace that germ in
the equivalence. This also transports the corresponding singular loci.
\end{proof}

\begin{definition}[The recorded-stage predicate]
\label{hir:recorded-stage-centres}
\lean{Hironaka.HironakaStageCenters.StageCentres}
\leanok
For a predicate $P(T,C)$ on an actual prefix tower and its next centre,
$\operatorname{StageCentres}(P,T)$ is true for the identity tower. For an
extended tower $T\pi$ with recorded centre $C$, it means
$\operatorname{StageCentres}(P,T)$ and $P(T,C)$. Thus it checks every actual
recorded centre and its own prefix projection and accumulated boundary.
No centre-containment evidence is built into the definition.
\end{definition}

\begin{theorem}[Every recorded centre outside the boundary]
\label{hir:stage-centres-outside-boundary}
\uses{hir:stage-cosupport-off-boundary,hir:stage-regularity-unramified}
\lean{Hironaka.HironakaStageCenters.stageCentresInStrictTransformOrBoundary_of_isGeoTestTower,Hironaka.HironakaStageCenters.stageCentresSingularOffBoundary_of_isGeoTestTower}
\leanok
For every actual multiple test tower of $(X,X,E_0,\mathcal I_W,1)$, each
recorded centre $C_j$ is contained in $W_j\cup|E_j|$. If in addition all
centres lie over
$W\setminus\operatorname{Reg}_d(W)\cup(W\cap|E_0|)$, then
\[
 C_j\setminus|E_j|\subset
 W_j\setminus\operatorname{Reg}_d(W_j).
\]
These are properties of every recorded step, not merely of the final support.
\end{theorem}
\begin{proof}
At a recorded step admissibility says $C_j\subset\operatorname{cosupp}
\mathfrak m_j$. Apply Lemma~\ref{hir:stage-cosupport-off-boundary}.
For the second assertion a point $c\in C_j\setminus|E_j|$ is in $W_j$
and is unramified for the prefix projection. Its image lies in the prescribed
original bad set. It cannot lie in the initial boundary: outside the
accumulated divisor the image is outside the initial boundary, by the same
induction used in Lemma~\ref{hir:stage-cosupport-off-boundary}. Thus its image
is singular of dimension $d$. Lemma~\ref{hir:stage-regularity-unramified}
shows that $c$ is singular in the current strict transform. Induct on the
tower to obtain the stated property at every prefix.
\end{proof}

\begin{theorem}[Compact pure-dimensional resolution with recorded positions]
\label{hir:embedded-stage-positions-off-boundary}
\uses{hir:stage-centres-outside-boundary,hir:embedded-resolution-pure-dimensional}
\lean{Hironaka.HironakaStageCenters.exists_tower_embeddedResolution_pure_stage_positions}
\leanok
A compact pure-dimensional analytic set $W$ of dimension $d<n$ has an actual
finite test tower, admissible over its singular set, with smooth strict
transform of dimension $d$ having normal crossings with the accumulated
divisor. Every recorded centre belongs to the current strict transform or
the current divisor, and its part outside that divisor is in the current
singular locus.
\end{theorem}
\begin{proof}
Initialize the actual vanishing ideal with mark one and empty boundary.
Cartan coherence and the pure-dimensional empty-interior theorem put it in
the domain. Its compact cosupport equals $W$. The canonical compact-cosupport
driver produces a finite actual resolution tower. Truncate it before its
first centre meeting the regular open set, using the existing identity-principle
construction. The resulting actual tower has smooth normal-crossings strict
transform, and all its centres avoid the preimage of that regular open set.
Thus they lie over the original singular set, which gives admissibility with
that bad locus. Apply Theorem~\ref{hir:stage-centres-outside-boundary} to the
recorded tower. No property of an exceptional centre inside the divisor is
silently inferred here.
\end{proof}

\begin{lemma}[A fresh canonical maximum detects a nonsmooth current set]
\label{hir:fresh-current-centre-singular}
\uses{hir:canonical-centres,hir:embedded-resolution-pure-dimensional}
\lean{Hironaka.HironakaStageCenters.fresh_maxLocus_subset_singular_of_not_isClosedSubmanifold,Hironaka.HironakaStageCenters.exists_fresh_admissibleCentre_subset_singular}
\leanok
Let $W$ be a compact analytic set of pure dimension $d<n$. If $W$ is not a
closed smooth submanifold, the actual canonical invariant, initialized on
$(X,X,[],\mathcal I_W,1)$, has a nonempty smooth admissible maximum locus
$C$ contained in the singular locus of the present $W$.
\end{lemma}
\begin{proof}
Cartan coherence and empty interior put the actual vanishing ideal in the
domain of the canonical invariant. Its initial cosupport is $W$. Upper
semicontinuity on the compact nonempty set $W$ produces a maximum, and the
admissibility theorem produces a closed smooth centre $C\subset W$.
Suppose that $C$ meets the regular locus. On the open complement of the
singular locus, regularity has constant dimension $d$; restriction to that
open set and invariance under local biholomorphisms make all its cosupport
points have the same invariant. Since one belongs to the global maximum,
all do. Hence $C$ and $W$ agree on that open set. The regular locus is dense
in $W$, and $C$ is closed, so $W\subset C$. The analytic filling argument
in the smooth manifold $C$, with the pure-dimensional density hypothesis,
shows that $W$ and $C$ have equal germs at every point of $W$. Therefore
$W$ is a closed smooth submanifold, contrary to hypothesis. Thus $C$
contains no regular point of $W$, proving $C\subset\operatorname{Sing}W$.
No premise asserts this centre containment: it is produced from the actual
canonical maximum and the geometry of its initial cosupport.
\end{proof}

\begin{theorem}[A genuine singular blowup and the next current ideal]
\label{hir:fresh-current-singular-blowup}
\uses{hir:fresh-current-centre-singular,hir:geo-transform}
\lean{Hironaka.HironakaStageCenters.exists_fresh_singular_blowup_subvariety}
\leanok
Let $V$ be a compact irreducible analytic subvariety of dimension $d<n$ which
is not smooth. There is a nonempty smooth centre $C\subset\operatorname{Sing}V$
and an actual blowup $\pi:Y\to X$ along $C$. The actual current strict
transform is the carrier of a compact irreducible analytic subvariety $V'$ of
dimension $d$. The new boundary $[\pi^{-1}C]$ has simple normal crossings,
and the freshly initialized marked ideal
$(Y,Y,[\pi^{-1}C],\mathcal I_{V'},1)$ belongs to the actual domain of the
canonical invariant.
\end{theorem}
\begin{proof}
Apply Lemma~\ref{hir:fresh-current-centre-singular} to $V$. The centre has
empty interior because it is contained in $V$ and $d<n$. Demailly's explicit
blowup construction therefore supplies the manifold $Y$ and map $\pi$.
Its exceptional image is contained in $C\subset\operatorname{Sing}V$.
Choose a regular point of $V$, which exists for an irreducible analytic
subvariety. It is outside that image, so $V$ is not contained in the
exceptional image. The strict-transform irreducibility and dimension theorem
then constructs $V'$ with the stated carrier and dimension. The strict
transform is closed inside $\pi^{-1}V$, which is compact by properness of
the actual blowup; hence $V'$ is compact.

The centre is admissible for $(X,X,[],\mathcal I_V,1)$, so its actual marked
transform is well formed and its boundary $[\pi^{-1}C]$ is SNC. The present
subvariety $V'$ is analytic of dimension $d<n$, so its actual vanishing ideal
is coherent and nonzero at every point. These properties, together with the
SNC boundary transverse to the whole ambient manifold, construct the domain
membership of the freshly initialized ideal. This statement constructs one
honest geometric step and its continuation data; it does not assert a bound
on repeated fresh initialization.
\end{proof}

\paragraph{The remaining literal producer.}
None of the preceding results closes the obligation on centres contained in
$|E_j|$. It is insufficient to identify the controlled cosupport with the
strict transform. For example, the reduced curve ideal
$(y^2-x^3,z)$ in $\mathbb C^3$, blown up at the origin, has controlled ideal
$(x(y'^2-x),z')$ in the $x$-chart, whereas its strict transform has ideal
$(y'^2-x,z')$. The controlled ideal has an additional exceptional branch
$\{x=z'=0\}$, and extracting divisorial common factors does not remove it.
This calculation illustrates the missing implication; it is not a claimed
counterexample to the actual maximum centre.

A dynamic-current-ideal route can apply the fresh-centre producer to each
actual current analytic strict transform, with the actual accumulated
boundary. To turn this into finite embedded resolution one must prove an
invariant of these freshly initialized current ideals that strictly
decreases under those steps, remains compatible with restriction and boundary
history, and stops precisely at the smooth normal-crossings locus. The
existing decrease theorem concerns controlled transforms of one initial
marked ideal and does not prove this reinitialization statement.
The second alternative, for smooth current $W_j$, further needs centres
contained in $W_j\cap|E_j|$; initial marked-ideal principalization by itself
can instead choose the whole smooth $W_j$. Both obligations remain explicit.
